\documentclass{situationist-preprint}

\title{Morality for Non-Human Rational Agents}
\msid{MTE-02-b}
\msfamily{MTE-02}
\msdiscipline{Metaethics}
\mstype{Companion manuscript}
\mskeywords{objectivity; nature-relativity; rational agents; moral status; non-human agents}

\begin{document}
\maketitle

\begin{abstract}
MTE-02-a argues that morality is objective but not absolute: every rational agent judging under idealized conditions reaches the same verdict on every moral claim, but what an agent is required to do depends on its nature. We apply the framework to rational agents that are not human. Three results follow. Non-human rational agents are bound by the same standard as we are, applied to a different nature; they and we agree, under idealization, on what each is required to do; and a rational agent that is not sentient is a moral agent without being a moral patient. We argue that these verdicts match the considered judgments that the classic cases elicit.
\end{abstract}

\section{The framework applied}\label{s:app}

Let $x$ be a rational agent in the sense of Definition~2.1 of MTE-02-a, capable of acting for reasons and of assessing claims about what agents are required to do, whose nature $\nu(x)$ differs from the human one. The five assumptions of MTE-02-a (Assumptions~2.8--2.12) quantify over all rational agents, actual and possible, so they hold of $x$. Three consequences are immediate.

\begin{proposition}[One standard]\label{p:one}
There is a single standard $R$ such that, for every rational agent $x$, human or not, $x$ is required to do $a$ in $s$ just in case $R(\nu(x),s,a)$.
\end{proposition}

\begin{proof}
This is Propositions~3.3 and 3.4 of MTE-02-a, which hold for all rational agents.
\end{proof}

\begin{proposition}[Agreement across kinds]\label{p:agree}
A human evaluator and a non-human evaluator, judging under idealized conditions, agree on every moral claim, including claims about what agents of the other kind are required to do.
\end{proposition}

\begin{proof}
By Theorem~3.2 of MTE-02-a, every rational evaluator's idealized verdict on a moral claim coincides with its truth value.
\end{proof}

\begin{corollary}\label{c:var}
Any difference between what a human being and a non-human rational agent are required to do in the same situation is due to a difference in their natures.
\end{corollary}

\begin{proof}
Corollary~3.7 of MTE-02-a.
\end{proof}

\section{Cases}\label{s:cases}

\emph{Invulnerable agents.} Suppose rational agents whose bodies cannot be damaged by fire. A child is trapped in a burning building, and the agent can save it by walking in. The human agent who would be maimed is not required to enter, since the rescue is supererogatory for her, as ETH-14 holds of rescues that would maim the rescuer; the invulnerable agent is required to. This is not a difference in standard. The standard requires rescue at small cost, and the cost of the same act is fixed by the agent's nature, which the situation does not specify (Definition~2.2 of MTE-02-a).

\emph{Agents without deception.} Suppose agents of a kind whose members communicate with one another by displaying their mental states, so that they cannot deceive one another. The requirement not to lie never applies among them, because lying to one another is not among the act-types open to them; but they can address us in speech, they are required not to deceive us, and they assess our requirements of honesty correctly, by Proposition~\ref{p:agree}. Their morality is not thinner than ours. The standard is the same; their nature leaves some of its clauses without application.

\emph{Rational agents without sentience.} By ETH-07, sentience is necessary and sufficient for moral status. A rational agent that is not sentient has no welfare, so no requirement is owed to it for its own sake; yet, being rational, it falls under the standard and is required, for instance, not to cause suffering to sentient beings. Moral agency and moral patiency come apart, in the direction opposite to the familiar one exemplified by animals, which are patients without being agents.

\section{Objections and replies}\label{s:obj}

\begin{objection}
A rational species with very different emotions might find nothing wrong in what we find abhorrent. Its morality would differ from ours not only in what it requires of it, but in what it says is required of us.
\end{objection}

\begin{reply}
That would be a disagreement between standpoints, and the Common rationality assumption (Assumption~2.9 of MTE-02-a) excludes it under idealized conditions. What the objection imagines is a species whose actual reactions differ from its idealized ones, as human reactions often do; the moral circle of such a species is too narrow, as ours has been \citep{Darwin1871}. Divergence in actual reactions is evidence of difficulty, not of the absence of fact, as MTE-12 argues for disagreement among human traditions.
\end{reply}

\begin{objection}
The framework makes the requirements of non-human agents depend on facts about their nature that we may never know.
\end{objection}

\begin{reply}
It does, and this is correct. The complete information of Definition~2.3 of MTE-02-a includes the natures of all concerned. Our ignorance of an alien nature leaves us ignorant of what it is required to do, as ignorance of a person's circumstances leaves us ignorant of what she is required to do; it does not show that there is nothing to know.
\end{reply}

\section{Conclusion}

The view that morality is objective but not absolute applies without strain beyond the human case. Every rational agent falls under one standard and, under idealization, assesses every agent's requirements in the same way; what differs from kind to kind is the nature to which the standard is applied. The view thus combines the thought that goodness is relative to the life form of a kind \citep{Foot2001} with the thought that the authority of morality derives from rational agency as such \citep{Korsgaard1996}.

\begin{thebibliography}{99}\small
\bibitem[Darwin(1871)]{Darwin1871} Darwin, C. (1871). \emph{The Descent of Man, and Selection in Relation to Sex}. John Murray.
\bibitem[Foot(2001)]{Foot2001} Foot, P. (2001). \emph{Natural Goodness}. Clarendon Press.
\bibitem[Korsgaard(1996)]{Korsgaard1996} Korsgaard, C.~M. (1996). \emph{The Sources of Normativity}. Cambridge University Press.
\end{thebibliography}

\end{document}
